\documentclass[11pt]{article}
\usepackage{explorer}
\newcommand{\ExplorerLanguage}{FR}
\newcommand{\ExplorerNumber}{08}
\newcommand{\ExplorerKind}{Leçon et méthode}
\begin{document}
\ExplorerTitle{Terminale · probabilités et suites}{Attendre le premier succès}
\ExplorerSection{Avant de commencer}
Épreuves indépendantes, suites géométriques et sommes.\par
\ExplorerSection{Idées et méthode}
On répète des essais indépendants ayant chacun une probabilité p de succès, avec \(0 < p \leq  1\). T est le numéro du premier succès, en commençant à 1.\par
L’événement \(T = k\) impose \(k-1\) échecs puis un succès : \(P(T=k) = (1-p\))ᵏ⁻¹p. De même, \(P(T>n) = (1-p)^{n}\) pour \(n \geq  0\).\par
L’espérance vaut \(E(T)=\frac{1}{p}\). On peut la retrouver par la somme des probabilités de queue : \(E(T)=\sum  P(T>n)\), pour \(n \geq  0\), donc par la somme géométrique. Il s’agit d’une moyenne, pas d’une date garantie de succès.\par
\ExplorerSection{Exemple expliqué}
Pour un dé équilibré et le succès « obtenir 6 », \(p=\frac{1}{6}\). L’attente moyenne est 6 lancers, mais la probabilité d’attendre plus de 6 lancers est \((\frac{5}{6})^{6}\), soit environ 0,335.\par
\ExplorerSection{À explorer}
Pour \(p=\frac{1}{4}\), comparer \(P(T>n)\) pour \(n=0,1\),…,10. Relier la diminution au facteur géométrique \(\frac{3}{4}\).\par
\end{document}
